% runtime_http_contract.tex —— 生产 HTTP/GUI 图分析与化简组合契约

\section{生产 HTTP 与 GUI 组合契约}\label{sec:gr-runtime-http}

\subsection{会话快照与并发口径}

生产路由实现在 GUI server 测试可执行的源文件中。拓扑分析读取会话内常驻的不可变
\texttt{PowerSystemGraph}；网络化简读取不可变系统快照并创建请求局部的 graph/system 副本。两条路由
都用会话级 \fld{busy} 排斥其他分析：已有任务时返回 HTTP 409 busy-conflict 错误。未加载系统、
JSON 解析失败或执行异常返回
HTTP 400。返回 200 只表示路由完成，不是 PF/OPF 或语义等价证书。

\subsection{POST /api/session/topology}

该路由忽略请求体，不接受算法选项，直接对 resident graph 调用
\srcpath{analyze\_topology}。顶层字段为：

\begin{longtable}{@{}L{4.4cm}L{4.2cm}L{5.6cm}@{}}
\toprule
字段 & 索引/单位 & 解释\\
\midrule
\endfirsthead
\toprule
字段 & 索引/单位 & 解释\\
\midrule
\endhead
\fld{is\_connected/is\_radial} & bool & 整个在运混合无向图的结论\\
\fld{cycle\_count} & count & $E_{on}-N_{on}+p$，不返回圈的边列表\\
\fld{n\_ac\_islands/n\_dc\_islands} & count & DC 只计纯 DC 岛；混合岛归 AC\\
\fld{all\_islands\_valid} & bool & 参考节点与孤立负荷综合标志\\
\fld{n\_buses/n\_branches} & graph storage count & \texttt{node\_count/edge\_count}，包含图中停运对象，不是 live count\\
\fld{islands[]} & stable bus IDs & 含 domain、status、AC/DC 母线列表和参考标志\\
\fld{ac/dc\_bus\_island[]} & authored positions & 按系统 AC/DC bus vector 顺序，停运/未归岛为 $-1$\\
\fld{cut\_vertices[]} & domain+stable ID & \texttt{bus/domain/pos}；推荐消费者使用\\
\fld{cut\_vertex\_bus\_ids[]} & flat stable IDs & 兼容输出，同号 AC/DC 时有歧义\\
\fld{ac/dc\_cut\_vertex[]} & authored positions & Canvas 布尔掩码\\
\fld{bridges[]} & stable endpoints+positions & 每端独立给出 domain、bus ID 和 Canvas position\\
\fld{diagnostics[]} & mixed legacy IDs & code/message/related buses/branches；related buses 未带域\\
\bottomrule
\end{longtable}

岛状态字符串为 \texttt{Valid/NoSlack/NoDCVoltageRef/IsolatedLoad/Empty}。带负荷、无本地源且
无任何关联边的单点孤点赋 \texttt{IsolatedLoad}（仅因支路停运而孤立、但仍带关联边的负荷母线归
\texttt{NoSlack}），优先于缺参考状态，并同步产生 GraphIsolatedLoad diagnostic 与
\fld{all\_islands\_valid=false}；\texttt{Empty} 仅是防御性分支。桥列表不暴露 graph edge position
或源组件 index，而是为 GUI
重建端点和类别；VSC 桥可有不同的 \fld{from\_domain/to\_domain}，legacy \fld{domain} 仅当两端均
为 DC 时写 DC。

\subsection{POST /api/session/network\_reduction 请求}

请求体为 JSON object，当前五个布尔选项如下：

\begin{center}\small
\begin{tabular}{@{}L{5.7cm}L{2.0cm}L{6.5cm}@{}}
\toprule
字段 & HTTP 默认 & 实际作用\\
\midrule
\fld{enable\_switch\_contraction} & true & 执行闭合开关/断路器收缩阶段\\
\fld{enable\_series\_reduction} & true & 单轮 classify/plan/apply series\\
\fld{enable\_pendant\_reduction} & false & 单轮悬垂负荷折叠\\
\fld{enable\_kron\_reduction} & false & 只标记响应 stage enabled；Kron 不执行\\
\fld{contract\_zero\_impedance\_lines} & false & 收缩阶段是否额外合并低阻 AC/DC line\\
\bottomrule
\end{tabular}
\end{center}

HTTP 与 C++ 默认的关键差异是 \fld{contract\_zero\_impedance\_lines}：C++ 为 true，而路由为
false，防止用户遗漏
$r/x$ 参数的普通线路被静默并母线。GUI 发出的请求不含
\fld{enable\_kron\_reduction}，所以 GUI 下该值始终为 false。即使为 false，路由仍执行一次候选分类并
返回 \fld{n\_candidates}；\fld{enabled} 表示请求意图，不是“候选扫描是否发生”。

\subsection{网络化简实际流水线}

可达执行顺序固定为：
\begin{equation}
G_0\xrightarrow[\text{可选低阻 line}]{\text{switch/breaker contraction}}
G_1\xrightarrow{\text{one series plan}}G_2
\xrightarrow{\text{one optional pendant plan}}G_3.
\end{equation}
每阶段都在前一阶段的新 graph/system 上重新 classify/plan。series 不是迭代到不动点；Kron 只在
$G_3$ 上统计 passive candidates，因为矩阵 Kron 不生成 rich system，故不会改变 after count 或
\fld{reduced\_system}。该路由不调用 PF/OPF，也不执行电压/支路结果恢复。

路由用 \srcpath{compose\_reduction\_mappings} 逐阶段复合公共 C++ mapping，并从合成后的
domain-qualified bus certificate 生成 authored bus $\to$ 当前代表展示；支路映射保留在 C++ 证书中，
当前 HTTP 响应不单独序列化该表。

\subsection{网络化简响应}

\begin{longtable}{@{}L{4.4cm}L{4.4cm}L{5.4cm}@{}}
\toprule
字段组 & 空间/单位 & 解释\\
\midrule
\endfirsthead
\toprule
字段组 & 空间/单位 & 解释\\
\midrule
\endhead
\fld{before/after} & live counts & 在运 bus/edge 数；与 topology 的 storage count 不同\\
\fld{n\_buses/branches\_eliminated} & count & before-after；不是 plan 的引用计数\\
\fld{reduction\_pct\_buses} & percent & 空图为 0\\
\fld{stages} & counts+fidelity strings & 每阶段 enabled、记录数和明确的精确/近似口径\\
\fld{formulation\_notes[]} & text & 四条模型说明；不是机器可验证证书\\
\fld{switch\_groups[]} & domain+stable IDs & super bus 与原成员\\
\fld{series\_records[]} & domain+stable IDs, pu & $r_{eq},x_{eq}$ 及精确性字符串\\
\fld{pendant\_records[]} & domain+stable IDs & 只声明 flat-voltage approximate，不给误差\\
\fld{ac/dc\_bus\_reduction[]} & authored positions+stable IDs & status、代表 ID 和代表 authored position\\
\fld{diagnostics[]} & legacy bare IDs & 汇总 contraction/series/pendant diagnostics\\
\fld{reduced\_system} & JSON object & compact 后的 rich-model 子集，可重新加载\\
\fld{reduced\_system\_schema} & string & 固定为 \texttt{hacdcpf\_system\_json\_v1}\\
\bottomrule
\end{longtable}

母线 status 取 \texttt{retained/merged/series\_eliminated/pendant\_eliminated}。代表 position 在原始
同域 bus vector 中查找；代表不存在时为 $-1$。响应没有统一 \texttt{schema}、\texttt{model\_scope}、
\texttt{validity} 或 approximation error 字段，消费端应保留 \fld{stages.fidelity} 和 diagnostics。

\subsection{紧凑 reduced\_system 的边界}

reducer 内部保留停运 audit ghosts；HTTP 在序列化前删除停运母线/支路与停运详细负荷，并验证所有
保留设备端点仍存在。收缩层会重映射 Transformer3W、LCC、EnergyRouter ports、MobileStorage 的
current/target bus 及其他已列 rich terminals；three-phase 子系统保留原结构，因为涉及相域母线的
收缩会在 C++ 入口 fail closed。若 compact 检出 LCC、EnergyRouter 或移动储能悬空端点，路由抛错，
不再静默过滤资产。投影 mapping/certificate/report 与 telemetry 仍不随 reduced system 导出。

E2E 固定检查 Transformer3W、LCC、EnergyRouter 与 three-phase 数据经化简、compact、reload 后数量
和稳定身份保留。虚拟耦合边只证明连接性，未构造这些设备的电气等值；HTTP Kron 仍是 identify-only，
pendant 仍是 flat-voltage loss approximation，不能把可重载性解释为任意分析的方程等价证书。

\subsection{GUI 与 topology\_window 的关系}

GUI 的拓扑面板调用 topology 路由并以 authored positions 高亮岛、桥和割点；网络化简面板显示
before/after、记录表和代表着色，并允许下载响应 JSON。它不会自动用 reduced system 替换当前会话。

\srcpath{POST /api/session/topology\_window} 只是 resident graph 空间索引上的 bbox/LOD 视窗查询：按
地理坐标筛选/聚合节点和边，不运行 graph reduction，也不产生恢复映射。它与本章两条分析路由共享
图数据来源，但不属于化简算法流水线。
